\documentclass[../main.tex]{subfiles}
\begin{document}

\section{Summary}

This thesis investigated the development of energy-efficient deep learning models
for capacitive image-based finger classification, with a particular emphasis on
Spiking Neural Networks (SNNs) as a viable alternative to conventional
convolutional and transformer-based approaches. The work was motivated by the
growing need for context-aware, low-power sensing in mobile and wearable devices,
where continuous finger-type recognition must operate within tight energy and
computational budgets.

Two distinct SNN-based architectures were proposed and evaluated on the
CapFingerId dataset. The first is a RANC-based SNN, designed as a fully spiking
architecture compatible with neuromorphic hardware platforms. The model employs
burst encoding to convert static capacitive images into spatiotemporal spike trains,
processes the encoded input through parallel overlapping segments via Tea layers,
and aggregates spiking activity using additive pooling for classification.
Experimental evaluation demonstrated that the proposed segmentation strategy ---
particularly the (\texttt{Elems=16}, \texttt{Overlap=8}) configuration --- achieves
the most favorable balance between recognition accuracy and model parameter
efficiency. The temporal encoding experiments further confirmed that moderate
timestep values ($T = 4$ to $T = 8$) provide meaningful performance gains, while
larger windows yield diminishing returns relative to their additional computational
cost. Overall, the RANC-based SNN achieved competitive accuracy compared to the
CNN baseline from the original CapFingerId study, while operating with
significantly fewer parameters and a neuromorphic-compatible design.

The second proposed architecture is a Spiking Vision Transformer (Spiking ViT),
built upon the SpikeFormer framework. This model integrates a Spiking Patch
Splitting (SPS) module for feature embedding and a Spiking Self-Attention (SSA)
mechanism to capture global spatial dependencies within the spike domain. By
adapting the encoder depth and embedding dimensions to suit the low-resolution
nature of capacitive data, the proposed Spiking ViT demonstrated consistent
improvement over the CapFingerId baseline across all evaluated classification tasks.
Notably, the best-performing configuration ($L = 8$, $D = 768$, $T = 4$) achieved
a Thumb L/R accuracy of $0.95$ and a Thumb/Index accuracy of $0.90$,
representing approximately $6\%$ improvement over the baseline model. The
experimental analysis showed that deeper encoder structures and larger embedding
dimensions generally improve performance, though gains saturate at high model
capacities, reinforcing the importance of selecting a balanced architecture.

Together, the results validate that SNN-based architectures --- particularly those
incorporating transformer-level attention mechanisms --- are capable of efficient and
accurate finger-type recognition from capacitive images. The event-driven and
sparse computation inherent to SNNs makes these models well-suited for deployment
on edge devices where energy constraints are critical.

Despite these contributions, the work has several limitations that should be
acknowledged. First, the actual deployment of the proposed RANC-based SNN on
FPGA hardware was beyond the scope of this thesis; the energy efficiency gains from
neuromorphic hardware therefore remain to be empirically verified. Second, the
Spiking ViT adopts a hybrid CNN-SNN design in its SPS module, which introduces
some non-spiking operations and partially limits the full spike-driven computation
ideal for neuromorphic hardware. Third, the study focused exclusively on tap
interactions from the CapFingerId dataset; generalization to dynamic gestures such
as swipes and scrolls was not evaluated.

\section{Suggestions for Future Work}

Several directions are proposed to extend and strengthen the findings of this
thesis.

The most immediate next step is the hardware deployment of the RANC-based SNN
on a physical neuromorphic platform, such as an FPGA running the RANC
architecture. This would enable direct measurement of energy-per-inference and
latency, providing concrete evidence of the power efficiency advantages that are
claimed theoretically in this work.

A second direction concerns the development of a fully spike-driven Vision
Transformer, eliminating the convolutional components in the SPS module and
replacing them with pure spiking operations. This would make the Spiking ViT
more compatible with neuromorphic hardware and further reduce power consumption
during inference.

Third, future work could explore the extension of both architectures to dynamic
gesture inputs --- including dragging and scrolling --- which introduce additional
temporal variability not addressed in this study. Encoding motion-dependent spike
patterns could further improve context awareness in real-world interaction
scenarios.

Finally, a systematic energy benchmarking study comparing the proposed SNN
models against CNN and standard ViT baselines in terms of millijoules per
inference ($\text{mJ}/\text{inference}$) would significantly strengthen the
practical case for neuromorphic finger recognition, closing an important gap
identified in the current literature.

\end{document}